% Image credits for the photographs and AI illustrations of Book 2 (High School), Hindi edition.
% Machine-readable license records live in images/book2/CREDITS.md; keep this
% file in sync with frontmatter/image-credits-book2.tex (one \item per item there).
% Names of people, institutions, collections, works, datasets and licences are
% kept exactly as in the English credits (Latin script); only the frame and the
% descriptions are translated (Book 1 hi convention).
\chapter*{चित्रों का श्रेय}

इस पुस्तक के रेखाचित्र, मानचित्र और चार्ट मौलिक हैं; प्रकाशित आँकड़ों से बनाए गए मानचित्र और
चार्ट अपना स्रोत चित्र-परिचय में या इस सूची में बताते हैं। फ़ोटो उनके अपने-अपने मुक्त लाइसेंसों के
अंतर्गत, उनके लेखकों के प्रति आभार के साथ, उपयोग किए गए हैं; हर फ़ोटो के पूरे स्रोत-लिंक और
लाइसेंस के पते पुस्तक के भंडार में \texttt{images/book2/CREDITS.md} में दिए गए हैं।

\bigskip
फ़ोटो और मौलिक रेखाचित्रों के साथ-साथ, कुछ चित्र एआई-जनित चित्रण हैं, जो पुस्तक के
संपादकों के लिखे निर्देशों से OpenAI के चित्र-निर्माण द्वारा बनाए गए, और शामिल करने से पहले उनकी
वैज्ञानिक शुद्धता जाँची गई। वे केवल सामान्य दृश्य दिखाते हैं, कभी कोई वास्तविक स्थान, वास्तविक व्यक्ति
या वास्तविक नमूना नहीं। हर जनित चित्र का पूरा निर्देश भंडार में
\texttt{images/book2/ai/PROMPTS.md} में दर्ज है।

\bigskip
\noindent फ़ोटो और आँकड़े, अध्याय के अनुसार:
\begin{itemize}\small\setlength{\emergencystretch}{2em}
  \item अध्याय~1 (पृथ्वी को मापना): Apollo~17 से देखी गई पृथ्वी; फ़ोटो: NASA, सार्वजनिक डोमेन (Wikimedia Commons)।
  \item अध्याय~1: एक पुनरावर्ती वृत्त, J.-B. Delambre की \emph{Base du système métrique décimal}, खंड~2 (1807) का एक चित्र-पट; सार्वजनिक डोमेन (Wikimedia Commons, BnF Gallica से)।
  \item अध्याय~1: WGS~84 दीर्घवृत्तज पर परिकलित अंश-लंबाइयाँ (EPSG dataset)।
  \item अध्याय~2 (एक वासयोग्य ग्रह): Mariner~10 से शुक्र (PIA23791), NASA/JPL-Caltech, सार्वजनिक डोमेन; Curiosity रोवर से Mount Sharp (PIA16769), NASA/JPL-Caltech/MSSS, सार्वजनिक डोमेन (Wikimedia Commons)।
  \item अध्याय~2: ग्रहों के आँकड़े NASA fact sheets से; जल की प्रावस्था-सीमाएँ IAPWS releases SR1-86 और R14-08 से; वासयोग्य क्षेत्र की सीमाएँ Kopparapu और अन्य (2013) से; Proxima~b, NASA Exoplanet Archive से।
  \item अध्याय~3 (पृथ्वी का ऊर्जा बजट): आर्कटिक समुद्री बर्फ़ पर पिघले जल के ताल (ICESCAPE mission, 2011); फ़ोटो: NASA Goddard Space Flight Center, सार्वजनिक डोमेन (Wikimedia Commons)।
  \item अध्याय~3: ऊर्जा-बजट के मान IPCC Sixth Assessment Report, Working Group~I, Figure~7.2 से; एल्बिडो University of Colorado Boulder (MOSAiC) और National Snow and Ice Data Center से; ऊष्मा-परिवहन Trenberth और Caron (2001) से; शुद्ध प्राथमिक उत्पादन Field और अन्य (1998) से; ऊष्मरासायनिक आँकड़े NIST Chemistry WebBook से।
  \item अध्याय~4 (बादल क्यों बनते हैं): साफ़ मौसम के कपासी मेघ, फ़ोटो: Ralph F. Kresge, NOAA; एक कपासी-वर्षी मेघ की निहाई, NOAA National Severe Storms Laboratory; दोनों सार्वजनिक डोमेन (Wikimedia Commons)।
  \item अध्याय~4: संतृप्ति IAPWS releases से परिकलित; बूँदों के आकार NASA Jet Propulsion Laboratory से।
  \item अध्याय~5 (तूफ़ान और उष्णकटिबंधीय चक्रवात): International Space Station से Hurricane Isabel, फ़ोटो: Mike Trenchard, NASA Johnson Space Center; Kansas में एक टॉरनेडो, फ़ोटो: Sean Waugh, NOAA National Severe Storms Laboratory; तड़ित, फ़ोटो: C.~Clark, NOAA National Severe Storms Laboratory; सभी सार्वजनिक डोमेन (Wikimedia Commons)।
  \item अध्याय~5: चक्रवातों के पथ IBTrACS version~4r01 से (Gahtan और अन्य, 2024; Knapp और अन्य, 2010; NOAA National Centers for Environmental Information); हरिकेन श्रेणियाँ NOAA National Hurricane Center से; हरिकेन की ऊर्जा और बनने की परिस्थितियाँ NOAA Atlantic Oceanographic and Meteorological Laboratory से; तड़ित-दर Christian और अन्य (2003) से; टॉरनेडो के तथ्य NOAA Storm Prediction Center से।
  \item अध्याय~6 (परतदार महासागर): एक नमूना-संग्राहक रोज़ेट, फ़ोटो: Pablo Clemente-Colon, NOAA Photo Library; एक अनुसंधान जहाज़ से छोड़ा गया एक प्रोफ़ाइलिंग प्लव, NOAA Global Ocean Monitoring and Observing Program; महासागरीय क्लोरोफ़िल 2002--2010, Jesse Allen और Robert Simmon, NASA Earth Observatory; सभी सार्वजनिक डोमेन (Wikimedia Commons)।
  \item अध्याय~6: महासागरीय प्रोफ़ाइल World Ocean Atlas 2023 से (Reagan और अन्य, 2024; NOAA National Centers for Environmental Information); समुद्री जल का घनत्व अंतरराष्ट्रीय अवस्था-समीकरण से (Fofonoff और Millard, 1983); महासागर का आयतन Charette और Smith (2010) से; प्लवों के तथ्य Argo Program Office से; उत्प्रवाह NOAA National Ocean Service और Kessler (2006) से।
  \item अध्याय~7 (मृदा, धीरे-धीरे बनने वाला संसाधन): Kerala में नाइस पर एक लैटेराइट प्रोफ़ाइल, फ़ोटो: Werner Schellmann, CC~BY-SA~2.5; मध्य Madagascar में सैप्रोलाइट, फ़ोटो: Diorit, CC~BY-SA~3.0; एक मृदा-प्रोफ़ाइल, USDA Natural Resources Conservation Service, सार्वजनिक डोमेन (सभी Wikimedia Commons)।
  \item अध्याय~7: मृदा-उत्पादन और अपरदन की दरें Montgomery (2007) से; खनिज-सूत्र Mineralogy Database (webmineral.com) से; सैप्रोलाइट की मोटाई USGS Ground Water Atlas of the United States से।
  \item अध्याय~8 (अपवाह द्रोणियाँ, नदियाँ और बाढ़): जलधारा का प्रवाह मापता एक जलविज्ञानी और एक मापन-केंद्र, U.S. Geological Survey, सार्वजनिक डोमेन (Wikimedia Commons)।
  \item अध्याय~8: Washington के निकट Potomac River के वार्षिक शिखर और दैनिक विसर्जन, U.S. Geological Survey National Water Information System, केंद्र 01646500; Mississippi का विसर्जन और द्रोणी Kammerer (1990, USGS) से; Amazon का विसर्जन Callède और अन्य (2010) से; बाढ़ के जल में सुरक्षा NOAA National Weather Service से।
  \item अध्याय~9 (मीठे जल के संसाधन): 2000 और 2018 में Aral Sea, NASA Earth Observatory (MODIS/Terra); Sublette, Kansas के निकट सिंचित खेत, NASA (ASTER/Terra), सार्वजनिक डोमेन (Wikimedia Commons)।
  \item अध्याय~9: विश्व के जल-भंडार Shiklomanov (1993) के अनुसार, U.S. Geological Survey Water Science School; वैश्विक जल-चक्र Trenberth और अन्य (2007) के अनुसार, जैसा Trenberth और Asrar (2014) में छपा; न्यूबियाई भूजल की आयु Sturchio और अन्य (2004) से; High Plains जलभृत के परिवर्तन McGuire (2017, USGS) से; Aral Sea का बजट Gaybullaev और अन्य (2012) और Micklin (2007) से।
  \item अध्याय~10 (महाद्वीपीय विस्थापन से समुद्री तल प्रसरण तक): वेधन-जहाज़ \emph{JOIDES Resolution}, Ocean Drilling Program, Texas A\&M University, सार्वजनिक डोमेन (Wikimedia Commons)।
  \item अध्याय~10: Cande और Kent (1995) का ध्रुवता काल-मापक्रम; Deep Sea Drilling Project (1970) के Leg~3 वेधन के परिणाम; महासागर-तल की आयु Seton और अन्य (2020) से, EarthByte, लाइसेंस CC~BY~4.0।
  \item अध्याय~11 (भूकंपीय तरंगें पृथ्वी का भीतरी भाग प्रकट करती हैं): Puerto Rico में एक अस्थायी भूकंपमापी लगाना; फ़ोटो: Thomas Pratt, U.S. Geological Survey, सार्वजनिक डोमेन (Wikimedia Commons)।
  \item अध्याय~11: Dziewonski और Anderson (1981) का Preliminary Reference Earth Model; P और S तरंगों के आनुभविक यात्रा-काल Australian Passive Seismic Server (AuSPASS) से।
  \item अध्याय~12 (स्थलमंडल और दुर्बलतामंडल): एक हिमशैल, NOAA, सार्वजनिक डोमेन (USGS Water Science School)।
  \item अध्याय~12: ऊष्मा-प्रवाह Davies और Davies (2010) से, लाइसेंस CC~BY~3.0; ऊष्मा-प्रवाह प्रांत Roy, Blackwell और Decker (1972) से; Hirschmann (2000) का प्रावार ठोसांक; भूपर्पटी की मोटाई और घनत्व Christensen और Mooney (1995) से; ऊँचाइयाँ ETOPO1 (NOAA) से।
  \item अध्याय~13 (प्लेटें और उनकी गतियाँ): Cascade Range में एक उपग्रह-स्थितिनिर्धारण केंद्र; फ़ोटो: Dave Sherrod, U.S. Geological Survey, सार्वजनिक डोमेन (Wikimedia Commons)। उत्तरी प्रशांत के तल का छायांकित उच्चावच, National Geophysical Data Center, NOAA, सार्वजनिक डोमेन (Wikimedia Commons)।
  \item अध्याय~13: Argus, Gordon और DeMets (2011) के प्लेट-घूर्णन ध्रुव NNR-MORVEL56; केंद्रों की स्थितियाँ Nevada Geodetic Laboratory से (Blewitt, Hammond और Kreemer, 2018); Clague और Dalrymple (1987) की ज्वालामुखियों की आयु, U.S. Geological Survey; Bird (2003) की प्लेट-सीमाएँ; स्थल Natural Earth से।
  \item अध्याय~14 (महासागरीय कटक): Galápagos Rift पर तकिया लावा, NOAA Okeanos Explorer Program, सार्वजनिक डोमेन; Juan de Fuca Ridge पर एक काला धूम्रक, University of Washington और NOAA, सार्वजनिक डोमेन; गैब्रो, फ़ोटो: Darla Sondrol (GeoDIL), CC0; बेसाल्ट में पेरिडोटाइट, फ़ोटो: James St.~John, CC~BY~2.0 (सभी Wikimedia Commons)।
  \item अध्याय~14: महासागर-तल की गहराइयाँ ETOPO1 (NOAA) से और आयु Seton और अन्य (2020) से, EarthByte, लाइसेंस CC~BY~4.0; Parsons और Sclater (1977) के गहराई-नियम; Hirschmann (2000) का प्रावार ठोसांक; Argus, Gordon और DeMets (2011) के घूर्णन ध्रुव NNR-MORVEL56।
  \item अध्याय~15 (क्षेपण क्षेत्र): ब्लूशिस्ट और एक्लोजाइट, फ़ोटो: James St.~John, CC~BY~2.0; Kesennuma, जापान में 2011 की सुनामी की क्षति, फ़ोटो: Dylan McCord, U.S. Navy, सार्वजनिक डोमेन (सभी Wikimedia Commons)।
  \item अध्याय~15: भूकंप USGS earthquake catalogue (ComCat) से; Bird (2003) की प्लेट-सीमाएँ और अभिसरण-दरें; महासागर-तल की आयु Seton और अन्य (2020) से, EarthByte, लाइसेंस CC~BY~4.0; Argus, Gordon और DeMets (2011) की प्लेट-चालें NNR-MORVEL56; स्थल Natural Earth से।
  \item अध्याय~16 (खनिज): क्वार्ट्ज़, ऑर्थोक्लेज़, बायोटाइट और ओलिवीन, फ़ोटो: Shannon Heinle, Darla Sondrol, Dexter Perkins और Nessa Eull (Geology Digital Image Library), CC0; पतली काट में बेसाल्ट, फ़ोटो: Chmee2, CC~BY-SA~3.0; पतली काट में एक ग्रेनाइट का क्वार्ट्ज़, फ़ोटो: Michael C.~Rygel, CC~BY-SA~3.0 (सभी Wikimedia Commons)।
  \item अध्याय~16: Rudnick और Gao (2014) के अनुसार महाद्वीपीय भूपर्पटी का संघटन; IUPAC/CIAAW के मानक परमाणु भार; खनिज-सूत्र और घनत्व webmineral.com से; Sierra Nevada की ग्रेनाइटी चट्टानों के घनत्व U.S. Geological Survey से (Open-File Report 91-570-A)।
  \item अध्याय~17 (आग्नेय चट्टानें): Giant's Causeway के बेसाल्ट स्तंभ, फ़ोटो: Adam Bishop, CC~BY-SA~3.0; El Capitan, Yosemite, U.S. Geological Survey, CC0; बेसाल्ट, ऐंडेसाइट, रायोलाइट, डायोराइट, ग्रेनाइट, ऑब्सीडियन, एक पॉर्फ़िरी चट्टान और दो पतली काटें, फ़ोटो: Nessa Eull, Twyla Baker, Joseph H.~Hartman, Mike Davis, Dexter Perkins और Darla Sondrol (Geology Digital Image Library), CC0 (सभी Wikimedia Commons)।
  \item अध्याय~17: डाइकों और सिलों का शीतलन Delaney (1986) के अनुसार, U.S. Geological Survey Open-File Report 86-444; मैग्मा-संघटन और क्रिस्टलन का क्रम S.~A.~Nelson's physical geology notes और S.~Earle, \emph{Physical Geology} (CC~BY~4.0) से; Bowen की जीवनी Carnegie Science history of the Geophysical Laboratory से।
  \item अध्याय~18 (कोयले, तेल और गैस का भूविज्ञान): कार्बन-समृद्ध शेल में जीवाश्म पौधे, लिग्नाइट, बिटुमिनी कोयला और ऐंथ्रेसाइट, फ़ोटो: James St.~John, CC~BY~2.0; Romania में एक लिग्नाइट संस्तर, U.S. Geological Survey, सार्वजनिक डोमेन (सभी Wikimedia Commons)।
  \item अध्याय~18: तेल-खिड़की, कोयले का बनना और ट्रैप S.~Earle, \emph{Physical Geology} (CC~BY~4.0) से; कोयले में कार्बन की मात्रा U.S. Energy Information Administration से; पेट्रोलियम तंत्र Magoon और Schmoker (U.S. Geological Survey, 2000) के अनुसार; पश्चिमी मध्य अफ़्रीका की दरार और लवण द्रोणियाँ U.S. Geological Survey Fact Sheet 2023-3028 से।
  \item अध्याय~19 (महाद्वीपीय भूपर्पटी): Acasta gneiss, फ़ोटो: Pedroalexandrade, CC~BY-SA~3.0; Jack Hills के ज़र्कॉन, चित्र: A.~Cavosie और J.~Valley, सार्वजनिक डोमेन; Canadian Shield के नाइस और ग्रेनाइट, फ़ोटो: Chris Light, CC~BY-SA~4.0 (सभी Wikimedia Commons)।
  \item अध्याय~19: Rudnick और Gao (2014) के अनुसार भूपर्पटी का संघटन; Christensen और Mooney (1995), White, McKenzie और O'Nions (1992) और Karplus और अन्य (2019) के अनुसार भूपर्पटी की मोटाई; Cogley (1984) के अनुसार महाद्वीपों का विस्तार; महासागर-तल की आयु Seton और अन्य (2020) से, EarthByte, लाइसेंस CC~BY~4.0।
  \item अध्याय~20 (पर्वतों का निर्माण): अंतरिक्ष यान Columbia से Mount Everest, NASA, सार्वजनिक डोमेन; Kali Gandaki गॉर्ज में वलित चट्टान-परतें, फ़ोटो: Jean-Marie Hullot, CC~BY~2.0; Glarus Thrust, फ़ोटो: Christian Heine, CC~BY-SA~3.0; Semail ओफ़ियोलाइट का तकिया लावा, फ़ोटो: Lionel Duchoiselle, CC~BY-SA~3.0 (सभी Wikimedia Commons)।
  \item अध्याय~20: भारत और ल्हासा की पिछली स्थितियाँ Müller और अन्य (2019, 2022) के प्लेट-मॉडलों से, GPlates Web Service (EarthByte) से परिकलित; ऊँची भूमि ETOPO1 से; भूकंप USGS catalogue से; Bird (2003) की प्लेट-सीमाएँ; स्थल Natural Earth से।
  \item अध्याय~21 (कायांतरित चट्टानें): सूक्ष्मदर्शी में स्लेटी विदलन, फ़ोटो: Geoloshe, CC0; Syros का गार्नेट-अभ्रक शिस्ट, फ़ोटो: Graeme Churchard, CC~BY~2.0; कायनाइट शिस्ट और क्वार्ट्ज़--बायोटाइट नाइस, फ़ोटो: James St.~John, CC~BY~2.0 (सभी Wikimedia Commons)।
  \item अध्याय~21: कायनाइट, ऐंडालुसाइट और सिलिमेनाइट के स्थायित्व-क्षेत्र Holdaway और Mukhopadhyay (1993) से।
  \item अध्याय~22 (क्षेपण क्षेत्रों पर मैग्मा-क्रिया): 18~मई 1980 को Mount St.~Helens के उद्गार का धूम-स्तंभ, फ़ोटो: Donald A. Swanson, USGS, सार्वजनिक डोमेन; ग्रैनोडायोराइट, ऐंडेसाइट और Sierra Nevada, फ़ोटो: James St.~John, CC~BY~2.0 (सभी Wikimedia Commons)।
  \item अध्याय~22: भूकंप USGS catalogue से; ज्वालामुखी Smithsonian Global Volcanism Program से; Bird (2003) की प्लेट-सीमाएँ; स्थल Natural Earth से; Hirschmann (2000) का शुष्क पेरिडोटाइट ठोसांक, PREM घनत्वों के साथ।
  \item अध्याय~23 (पर्वतों का विनाश): International Space Station से Appalachian Mountains, अंतरिक्ष-यात्री फ़ोटो ISS033-E-22378, NASA, सार्वजनिक डोमेन; क्वार्ट्ज़-कंकड़ संगुटिकाश्म, फ़ोटो: James St.~John, CC~BY~2.0 (दोनों Wikimedia Commons)।
  \item अध्याय~23: Portenga और Bierman (2011) की अनाच्छादन दरें; Meade और Moody (2010) के Mississippi अवसाद-भार; Ahnert (1970) का क्षय-स्थिरांक; वार्डन पहाड़ियाँ, श्रृंखला का कल्पित क्षेत्र।
  \item अध्याय~24 (पृथ्वी की आंतरिक ऊष्मा): Nesjavellir भूतापीय बिजलीघर, आइसलैंड, फ़ोटो: Gretar Ívarsson, संपादन: Fir0002, सार्वजनिक डोमेन (Wikimedia Commons)।
  \item अध्याय~24: Davies और Davies (2010) का ऊष्मा-प्रवाह; रेडियोजनित ऊष्मा KamLAND geoneutrino measurement (2011) से और चारों समस्थानिकों के वर्तमान ऊष्मा-उत्सर्जन से; चट्टानों का ऊष्मा-उत्पादन Vilà और Fernández (2010) से।
  \item अध्याय~25 (रेडियोमितीय काल-निर्धारण और पृथ्वी की आयु): William Thomson, Lord Kelvin, 1906, फ़ोटो: Russell and Sons, London, Smithsonian Institution Libraries, कोई ज्ञात कॉपीराइट प्रतिबंध नहीं; Canyon Diablo लौह उल्कापिंड, फ़ोटो: Geoffrey Notkin, CC~BY-SA~2.5 (दोनों Wikimedia Commons)।
  \item अध्याय~25: Patterson (1956) के उल्कापिंडों के सीसा-समस्थानिक; IUGS और IUPAC (2022) द्वारा अनुशंसित रुबिडियम-87 का क्षय-स्थिरांक; अनुमानों का इतिहास Holmes (1913, 1927), Patterson (1956) और U.S. Geological Survey से; वार्डन पहाड़ियाँ, श्रृंखला का कल्पित क्षेत्र।
  \item अध्याय~26 (आरंभिक पृथ्वी और उसका वायुमंडल): पट्टित लौह संरचना, Dales Gorge, फ़ोटो: Graeme Churchard, CC~BY~2.0; Shark Bay में स्ट्रोमैटोलाइट, फ़ोटो: Paul Harrison, CC~BY-SA~3.0; Sedona के निकट लाल संस्तर, फ़ोटो: Wendy from Pennsylvania, CC~BY-SA~2.0 (सभी Wikimedia Commons)।
  \item अध्याय~26: Gough (1981) के अनुसार युवा सूर्य की दीप्ति; इयॉन की सीमाएँ और रंग International Chronostratigraphic Chart से।
  \item अध्याय~27 (अतीत की जलवायु पढ़ना): अंटार्कटिक बर्फ़ में वायु के बुलबुले, फ़ोटो: CSIRO Atmospheric Research, CC~BY~3.0; प्लवकी फ़ोरैमिनिफ़ेरा, फ़ोटो: Marci Robinson, USGS, सार्वजनिक डोमेन; पराग-कण, Dartmouth Electron Microscope Facility, सार्वजनिक डोमेन (सभी Wikimedia Commons)।
  \item अध्याय~27: EPICA Dome~C का तापमान (Jouzel और अन्य, 2007) और अंटार्कटिक बर्फ़ के कोरों की कार्बन डाइऑक्साइड (Bereiter और अन्य, 2015), NOAA Paleoclimatology के ज़रिए; LR04 stack (Lisiecki और Raymo, 2005); ETOPO1 उच्चावच (NOAA); NASA GISS (GISTEMP v4); संदर्भ काल IPCC Sixth Assessment Report, Working Group~I, chapter~2 से।
  \item अध्याय~28 (जलवायु परिवर्तन): Mauna Loa Observatory; फ़ोटो: Brian Vasel, NOAA, सार्वजनिक डोमेन (Wikimedia Commons)।
  \item अध्याय~28: 1997 में Mauna Loa पर Charles David Keeling और John Chin; फ़ोटो: John Miller, NOAA, सार्वजनिक डोमेन (Wikimedia Commons)।
  \item अध्याय~28: आँकड़े NOAA Global Monitoring Laboratory और Scripps Institution of Oceanography (कार्बन डाइऑक्साइड), NASA GISS (GISTEMP v4), NOAA Laboratory for Satellite Altimetry (समुद्र-तल), Bereiter और अन्य (2015), NOAA Paleoclimatology के ज़रिए (बर्फ़ के कोर), और IPCC Sixth Assessment Report, Working Group~I (प्रणोदन, आरोपण, प्रक्षेपण) से।
\end{itemize}
